% ============================================================================
% Abstract
% ============================================================================
\section*{Abstract}

This thesis presents the design, implementation, and experimental evaluation of Coroute --- a high-performance C++ library for building web applications based on modern C++20 features. The primary objective is to demonstrate that the combination of stackless coroutines, OS-native asynchronous I/O mechanisms, and algorithmically optimized routing can achieve performance comparable to leading industrial solutions while significantly simplifying the programming model.

The architecture of Coroute is founded on several key innovations. The execution model employs the \texttt{Task<T>} type, implementing lazy coroutines with support for detached execution and cooperative cancellation. This model eliminates the traditional drawbacks of callback-based asynchronous programming while preserving the efficiency of non-blocking I/O. The network layer abstracts platform-specific mechanisms --- IOCP for Windows and io\_uring for Linux --- behind a unified interface, enabling zero-copy operations and syscall batching. Routing utilizes a DFA-based algorithm, published in ``Matching Text from Start to Finish Against Multiple Regular Expressions'' \cite{stankov2024regex}, achieving O(N) time complexity relative to URL length, independent of the number of registered routes. Fourth, the architecture integrates a template view system and seamless Flutter support via Dart FFI, enabling the reuse of the library as a headless web engine for building native mobile and desktop applications.

The library supports multiple protocols on a single port via ALPN negotiation: HTTP/1.1 with keep-alive, HTTP/2 with HPACK compression and stream multiplexing, and WebSocket for bidirectional real-time communication. Type-safe parameter extraction leverages C++20 concepts for compile-time validation, eliminating an entire class of runtime errors. The template view system extends server-side rendering capabilities through integration with the Inja template engine.

Experimental evaluation on a system with AMD Ryzen 5 3600 (6 cores, 12 threads) and up to 1024 concurrent clients demonstrates throughput of 1,378,578 requests per second for simple HTTP responses, with median latency of 212 microseconds and 0\% error rate. Comparative analysis with Crow, Boost.Beast, and Drogon shows that Coroute achieves comparable or superior performance with a significantly simpler API.

The scientific contributions include: integration of DFA-based routing into a full-featured web framework; a coroutine execution model specifically designed for network applications; and demonstration of practical applicability through a production-scale example application.

\vspace{1cm}

\textbf{Keywords:} C++20, coroutines, web server, HTTP/2, WebSocket, IOCP, io\_uring, DFA routing, type safety, asynchronous I/O, Flutter, Dart FFI, Template Views, Inja
